% Mass-spring-damper diagram.
% Usage: \massSpringDamper           (default scale 1.0)
%        \massSpringDamper[0.7]      (custom scale)
\newcommand{\massSpringDamper}[1][1.0]{%
    \begin{circuitikz}[scale=#1, transform shape]
        % --- ground / floor ---
        \fill[pattern=north east lines] (0,-0.25) rectangle (6,0);
        \draw[thick] (0,0) -- (6,0);

        % --- fixed wall on the left ---
        \fill[pattern=north east lines] (0,0) rectangle (0.25,2.6);
        \draw[thick] (0.25,0) -- (0.25,2.6);

        % --- mass block ---
        \def\mx{3}   % left edge of mass
        \draw[thick, fill=brandred20] (\mx,0.0) rectangle (\mx+1.4,2.15);
        % \node at (\mx+0.7,1.0) {$m = \SI{1}{kg}$};
        \node at (\mx+0.7,1.0) {$m$};

        % --- spring (top) and damper (bottom) in parallel ---
        \draw (0.25,1.6) to[spring, l={$k$}] (\mx,1.6);
        % \draw (0.25,1.6) to[spring, l={$k=\SI{3}{N/m}$}] (\mx,1.6);
        \draw (0.25,0.5) to[damper, l={$c$}] (\mx,0.5);
        % \draw (0.25,0.5) to[damper, l={$c=\SI{2}{N.s/m}$}] (\mx,0.5);

        % --- applied force F(t) ---
        \draw[very thick, red, -{Stealth[length=3mm]}] (\mx+1.4,1.0) -- (\mx+2.6,1.0)
            node[midway, above, red] {$f(t)$};

        % --- displacement x(t) ---
        \draw[dashed] (\mx+0.7,2.15) -- (\mx+0.7,2.9);
        \draw[-{Stealth[length=2mm]}] (\mx+0.7,2.5) -- (\mx+1.7,2.5)
            node[midway, above] {$x(t)$};
    \end{circuitikz}%
}
